% Method Section - Enhanced with learnable approaches

\subsection{The Limitation of Relation-Agnostic Variance}

Recall from \S\ref{sec:background} that GP-KGE uncertainty averages entity variances: $U_{\text{GP}}(h, r, t) = \frac{1}{2}(\bar{\sigma}^2_h + \bar{\sigma}^2_t)$ where $\bar{\sigma}^2_e$ depends only on entity $e$'s learned parameters.
This formulation has a critical limitation:

\begin{proposition}[Relation-Agnostic]
\label{prop:relation_agnostic}
GP-KGE uncertainty (Eq.~\ref{eq:gp_uncertainty}) is relation-agnostic: for any entity $e$ and relations $r_1, r_2$,
\begin{equation}
    U_{\text{GP}}(e, r_1, \cdot) = U_{\text{GP}}(e, r_2, \cdot)
\end{equation}
\end{proposition}

\begin{proof}
From Eq.~\ref{eq:gp_uncertainty}, $U_{\text{GP}}(h, r, t) = \frac{1}{2}(\bar{\sigma}^2_h + \bar{\sigma}^2_t)$ where $\bar{\sigma}^2_e = \frac{1}{d}\sum_j \exp(\ell_{e,j})$ is the mean variance for entity $e$.
Since $\bar{\sigma}^2_e$ depends only on entity parameters $\ell_e$ and not on any relation $r$, we have $U_{\text{GP}}(h, r_1, t) = U_{\text{GP}}(h, r_2, t)$ for all $r_1, r_2 \in \mathcal{R}$.
\end{proof}

This limitation is fundamental to all entity-level uncertainty methods, including BEUrRE~\citep{chen2021probabilistic} and UKGE~\citep{chen2019embedding}.
Making variance relation-aware naively would require $|\mathcal{E}| \times |\mathcal{R}| \times d$ parameters---infeasible for typical KG sizes.
We address this through four increasingly sophisticated approaches.

\subsection{Approach 1: Coverage-Based Structural Uncertainty}

The simplest fix is an explicit structural signal that captures entity-relation co-occurrence.

\begin{definition}[Binary Coverage]
For entity $e$ and relation $r$:
\begin{equation}
    c(e, r) = \mathbb{1}\left[\exists (h', r, t') \in \mathcal{T} : h' = e \lor t' = e\right]
\end{equation}
\end{definition}

Coverage equals 1 if entity $e$ appears with relation $r$ in training, else 0.
Structural uncertainty is then:
\begin{equation}
    U_{\text{cov}}(h, r, t) = 2 - c(h, r) - c(t, r)
\end{equation}

\paragraph{Frequency-Based Alternative.}
A natural extension replaces binary coverage with normalized frequency:
\begin{equation}
    c_{\text{freq}}(e, r) = \frac{\text{count}(e, r)}{\max_{e'} \text{count}(e', r)}
\end{equation}
where $\text{count}(e, r)$ is the number of training triples containing entity $e$ with relation $r$.
This distinguishes ``seen once'' from ``seen 1000 times.''
We compare both variants in \S\ref{sec:experiments}; binary coverage performs comparably, suggesting the observation/non-observation distinction matters more than frequency.

\paragraph{CAGP: Linear Combination.}
Our \textbf{C}overage-\textbf{A}ugmented \textbf{G}aussian \textbf{P}rocess (CAGP) approach linearly mixes both signals:
\begin{equation}
    U_{\text{CAGP}}(h, r, t) = \alpha \cdot \tilde{U}_{\text{GP}}(h,r,t) + (1-\alpha) \cdot U_{\text{cov}}(h,r,t)
\end{equation}
where $\alpha \in [0,1]$ is learned and $\tilde{U}_{\text{GP}}$ is normalized to match the coverage scale:
\begin{equation}
    \tilde{U}_{\text{GP}}(h,r,t) = U_{\text{GP}}(h,r,t) \cdot \frac{\bar{U}_{\text{cov}}}{\bar{U}_{\text{GP}}}
\end{equation}
with $\bar{U}_{\text{cov}}$ and $\bar{U}_{\text{GP}}$ denoting means computed over training triples (fixed at inference time). This ensures both signals contribute equally when $\alpha = 0.5$.

\paragraph{Training $\alpha$.}
We parameterize $\alpha = \sigma(\alpha_{\text{logit}})$ where $\sigma$ is the sigmoid function and $\alpha_{\text{logit}} \in \mathbb{R}$ is a learnable scalar initialized to 0 (i.e., $\alpha_0 = 0.5$).
After training the base GP-KGE model, we freeze embeddings and optimize $\alpha_{\text{logit}}$ on a held-out calibration set (10\% of validation data) using binary cross-entropy between predicted uncertainty and OOD labels.
Learned values vary by dataset structure: $\alpha = 0.38$ (FB15k-237), $\alpha = 0.52$ (WN18RR), $\alpha = 0.41$ (YAGO3-10)---see Appendix~\ref{app:alpha} for analysis.

While effective, this approach has limitations:
\begin{itemize}
    \item Coverage is binary---it doesn't distinguish between entities seen once vs. 1000 times with a relation
    \item Global $\alpha$ assumes uniform signal quality across all queries
\end{itemize}

\subsection{Approach 2: Attention-Based Query-Specific Mixing}

Instead of a global $\alpha$, we learn query-specific mixing weights:
\begin{equation}
    \alpha(h, r, t) = \sigma\left(\text{MLP}\left([\vect{h}; \vect{r}; \vect{t}; U_{\text{GP}}; U_{\text{cov}}]\right)\right)
\end{equation}

The attention mechanism considers:
\begin{itemize}
    \item \textbf{Entity/relation embeddings}: semantic content of the query
    \item \textbf{Current uncertainty values}: meta-information about signal quality
\end{itemize}

This allows the model to dynamically weight signals---using GP more heavily for rare-entity queries and coverage more heavily for novel-context queries.

For additional expressiveness, we employ multi-head attention with $K=4$ heads:
\begin{equation}
    \alpha(h,r,t) = \sum_{i=1}^K w_i(h,r,t) \cdot \alpha_i(h,r,t)
\end{equation}
where each head $\alpha_i$ is a separate MLP and the head weights $w_i(h,r,t) = \text{softmax}(\text{MLP}_w([\vect{h}; \vect{r}; \vect{t}]))_i$ allow different heads to specialize in different query patterns.

\subsection{Approach 3: Relation-Conditioned Variance}

Rather than post-hoc combination, we directly learn relation-aware variance:
\begin{equation}
    \sigma^2(e, r) = \text{softplus}\left(\text{MLP}\left([\vect{e}; \vect{r}]\right)\right) + \epsilon
\end{equation}
where $\epsilon = 10^{-6}$ ensures numerical stability.

This addresses the fundamental limitation: variance now depends on the $(e, r)$ pair, not just $e$.
The network learns to output:
\begin{itemize}
    \item Low variance when entity frequently appears with relation (well-constrained)
    \item High variance when entity rarely/never appears with relation (uncertain)
\end{itemize}

Uncertainty for a triple becomes:
\begin{equation}
    U_{\text{RelCond}}(h, r, t) = \frac{\sigma^2(h, r) + \sigma^2(t, r)}{2}
\end{equation}

This is trained jointly with the KGE objective, with the variance network learning to predict uncertainty from the embedding space structure.

\subsection{Approach 4: GNN-Based Uncertainty Propagation (GNNUncertainty)}

The most expressive approach propagates uncertainty through the graph:
\begin{equation}
    \sigma^{(l+1)}_e = \text{Update}\left(\sigma^{(l)}_e, \text{Aggregate}\left(\{\sigma^{(l)}_n : n \in \mathcal{N}(e)\}\right)\right)
\end{equation}

Key insight: an entity connected to many uncertain entities should itself be more uncertain.
This captures higher-order structural patterns beyond simple coverage.

The \textbf{GNNUncertainty} architecture consists of:
\begin{enumerate}
    \item \textbf{Initial uncertainty}: $\sigma^{(0)}_e$ from GP variance + coverage
    \item \textbf{Message passing}: relation-specific attention over neighbor uncertainties (2 layers)
    \item \textbf{Readout}: relation-conditioned MLP for final uncertainty $U_{\text{GNN}}(h,r,t)$
\end{enumerate}

This is the most computationally expensive variant but captures graph-level uncertainty patterns.

\subsection{Complementarity of Signals}
\label{sec:complementarity}

We observe empirically that semantic (GP) and structural (coverage) uncertainty capture \emph{different} failure modes.
Two illustrative cases demonstrate why neither signal alone suffices:

\paragraph{Case 1: Coverage succeeds, GP fails.}
Entity $e$ appears in 1000 triples with relations $\{r_1, \ldots, r_{10}\}$, but never with $r_{11}$.
GP assigns low variance $\sigma^2_e$ (well-observed entity).
For query $(?, r_{11}, e)$: GP signals low uncertainty, but the entity-relation pair was never observed---coverage correctly signals high uncertainty.

\paragraph{Case 2: GP succeeds, Coverage fails.}
Entity $e'$ appears in exactly one triple $(h, r, e')$.
Coverage assigns $c(e', r) = 1$ (observed).
For query $(h', r, e')$ with corrupted head: coverage signals low uncertainty, but GP correctly signals high uncertainty due to large $\sigma^2_{e'}$ from sparse observations.

\paragraph{Interpretation.}
GP captures \emph{how well an entity is learned} (global observation frequency), while coverage captures \emph{whether a specific context was observed} (local relation-specific).
These are orthogonal axes: Table~\ref{tab:temporal} validates this empirically, showing GP excels on emerging entities (0.83 AUROC) while coverage excels on novel contexts (1.00 AUROC).
